\begin{theorem}\label{night:P3K.23}
\StatementRecord{night:P3K.23}{false as stated; correction proved}{S6\_06\_BLIND: GPT-6 Astra (OpenAI, ChatGPT Pro mode, effort Pro), recorded re-derivation in the checkpoint dated 18 November 2026}{\cref{night:P3F.22}}\par\noindent\textbf{Physical-line correction, 17 November 2026.} The physical fixed-edge formula is \(7A_M\), where \(A_M\) is the actual physical integer. The old \(-5A\) formula is refuted. All displayed quaternion, common-link, relative-group and oriented commutator maps are retained. The native right integer is \(n=6A_M=-6A_U\); K68--69 prove the line and gluing dictionary and the corrected examples. The following earlier assertion and its argument are retained as the located historical claim; their corrected current consequences are proved in K68--69.\medskip\par 
Consider one of the actual sphere families constructed above with a
semifree fixed curve \(F\simeq\mathbb P^1\), whose complex normal bundle
is \(\mathcal O(-1)\oplus\mathcal O(-1)\). Perform the toric flop of
this fixed edge in the periodic cusp model, retaining the actual
parameter clutch and all exterior gluings. Write the original positive
return's Jin index as
\[
             A=\frac{4k-\ell}{24},
\]
The positive return of the flopped family has oriented class
\[
                         -5A\pmod {28}                 \tag{K23.1}
\]
in the run's normalization of \(\Theta _7\). The flopped curve has two
isolated circle fixed points. Formula (K23.1) specifies its homotopy
sphere class; it does not assign a new Jin integer to that circle
action.

In particular, the four groups of primitive-source examples in
P3K.22 give the following flopped classes:


\begin{center}\small\begin{tabular}{c l l}\hline
\(d=1/h(G)\) & Original index \(A\) & Flopped class in \(\mathbb Z/28\) \\
\hline
\(8\) & \(-3\) & \(15\) \\
\(14\) & \(-6\) & \(2\) \\
\(20\) & \(-9\) & \(17\) \\
\(30\) & \(-14\) & \(14\) \\
\hline\end{tabular}\end{center}

The earlier degree-four OS47 example of index \(-4\) gives class
\(20\). Every original zero-index example remains class zero after
this fixed-edge flop. No formula for a flop at another toric edge is
asserted here.
\end{theorem}

\begin{proof}
\begin{enumerate}[label=\arabic*.]
\item \par\medskip\noindent\textbf{The actual two local products and the parameter clutch.}\enspace
Write the conifold link as
\[
 K=(S^3_a\times S^3_c)/S^1,\qquad
 (a,c)\sim(e^{i\phi}a,e^{-i\phi}c).
\]
With \(\mathbb H\) oriented by \((1,i,j,k)\), put
\[
 u=a_0+j a_1,\quad v=c_0+j\overline{c_1},\quad
 q=uv,\quad h=u i u^{-1},\quad k=v^{-1} i v.
                                                        \tag{K23.2}
\]
The first resolution has the actual product coordinates \((h,q)\)
on its real normal sphere bundle; the second has \((k,q)\).
To invert the first map, choose a unit representative \(u\) of \(h\)
and put \(v=u^{-1}q\). Changing the representative changes only the
displayed auxiliary-circle orbit. For the second choose \(v\) and
put \(u=qv^{-1}\). These inverse formulas are smooth in the standard
quotient charts. The common boundary map is consequently
\[
             \tau(h,q)=(q^{-1}h q,q).                  \tag{K23.3}
\]
The fixed sphere is unknotted in the smooth six-sphere:
Haefliger's Theorem 1.2 and Corollary 6.6, with \(n=2,q=4\),
give the trivial embedding class \cite{Haefliger1966}.
Its exterior is therefore the actual product
\(E=D^3_y\times S^3_q\), with its given boundary marking extended
by the normal framing. The map (K23.3) extends over this exterior:
\[
 T(y,q)=(q^{-1}yq,q),\qquad
 T^{-1}(y,q)=(q yq^{-1},q).                           \tag{K23.4}
\]
These are smooth also at \(y=0\). Thus replacing the local resolution
conjugates the exterior return by \(T\).

\item This last sentence includes the actual parameter family. For a
local Cox angular action of weights
\((-r-s,0,r,s)\), write
\[
 a_\theta=e^{-i(r+s)\theta/2},\qquad
 b_\theta=e^{i(r-s)\theta/2}.
\]
The first product action and second product action are respectively
\[
 (h,q)\longmapsto
 (a_\theta h a_\theta^{-1},a_\theta q b_\theta),\qquad
 (k,q)\longmapsto
 (b_\theta^{-1}k b_\theta,a_\theta q b_\theta).
\]
Substitution in (K23.3) proves exact equivariance for every \(\theta\).
At \(2\pi\), both scalar endpoints have the same central sign;
the displayed product actions close even when \(r+s\) is odd.
Therefore the comparison of the two local family products is the
same map \(\tau\), independent of the parameter angle. The period
reindexing acts on the periodic collection of common conifold cells.
That collection has trivial cell stabilizer; choose one quotient
representative and extend every chart change equivariantly to its
translates. Extracting the integer winding from each bounded
nonzero toric unit leaves a logarithm on this representative. Contract
its positive modulus and its remaining logarithm in both resolution
charts at once, transporting the truncation collars by the resulting
common conifold isotopy. The integer winding is exactly the Cox
action just displayed. This retains the actual exterior
identifications, including their bounded units. No null-homotopy of
an odd Cox loop, or additional disc-supported correction, is inserted.

\par\medskip\noindent\textbf{Exact relative generators and conjugation.}\enspace
Choose a smooth degree-one map
\(\zeta:(D^3,S^2)\to(S^3,1)\), constant on a boundary collar,
which is equivariant for quaternionic conjugation and commutes
with \(y\). Explicitly one can use
\[
 \zeta(y)=\cos\phi(|y|)-\sin\phi(|y|)\frac{y}{|y|},
\]
where \(\phi\) decreases from \(\pi\) to zero, has the smooth odd
radial expansion at zero, and is flat near the boundary. Define
\[
 L(y,q)=(y,\zeta(y)q),\qquad R(y,q)=(y,q\zeta(y)).
                                                        \tag{K23.5}
\]
They commute and extend over \(D^3\times D^4\) by the same quaternion
multiplications. Their induced six-sphere classes are therefore zero.
They represent the left and right generators of
\(\pi_3SO(4)=\mathbb Z^2\). Extension by the identity to the double
injects the relative pseudo-isotopy group into the doubled group;
Krylov's Proposition 1 and Lemma 3 give, in this dimension, the
relative group \(\Theta _7\oplus\mathbb Z^2\)
\cite{Krylov2004}. Pseudo-isotopy is sufficient here, since extending
the resulting collar gives the same oriented sphere gluing.

\item On the double write
\[
 {\cal A}(s,t)=(tst^{-1},t),\quad
 {\cal B}(s,t)=(s,sts^{-1}),\quad
 {\cal C}=[{\cal A},{\cal B}].
\]
Krylov's Theorems 2--3 establish centrality and order \(28\) of
\({\cal C}\), and the relation
\({\cal A}Y=Y{\cal A}{\cal B}\), where \(Y(s,t)=(s,st)\)
\cite{Krylov2001}. In the marked relative products
\[
       T={\cal A}^{-1},\qquad L=Y,\qquad
       R=Y{\cal B}^{-1}.
\]
The first identity below also follows by literal substitution in
(K23.4)--(K23.5); the second then follows from the displayed central
relation:
\[
 TLT^{-1}=R,\qquad
 TRT^{-1}=L^{-1}R^2{\cal C}.
\]
Consequently
\[
 T(L^mR^n{\cal C}^s)T^{-1}
       =L^{-n}R^{m+2n}{\cal C}^{s+n}.                \tag{K23.6}
\]
All maps on the left are relative to the boundary. The injection
just cited retains this equality in the relative group. Transport
of a supporting oriented disc shows that \(T\) fixes the central
disc-supported subgroup.

\par\medskip\noindent\textbf{The actual integer coefficient from the hopf quotient.}\enspace
The orientations in (K23.2) are essential. At \(a_1=x+iy=0\),
the derivatives of \(h\) are \((-2k,-2j)\), so \(h\) reverses the
complex curve orientation. At \(u=1\),
\[
 q=\Re c_0+i\Im c_0+j\Re c_1+k\Im c_1
\]
preserves the complex normal orientation. Thus the complex
orientation of the exterior is \(D^3_y\times S^3_q\).
For the second resolution the derivatives of \(k\) are \((2k,-2j)\);
its base preserves and its normal coordinate reverses the respective
complex orientations. Its exterior has the same product orientation.

\item The physical delta action is \(q\mapsto q e^{-i\theta}\), with
Hopf quotient \(z=q i q^{-1}\). At \(q=1\), its circle vector is
\(-i\), and the quotient derivatives in the \(j,k\) directions are
\(-2k,2j\). Hence
\[
 [S^3_q]=-[S^1_\delta][S^2_z],\qquad
 [E]=[S^1_\delta][D^3_y][S^2_z].                    \tag{K23.7}
\]
The closed conifold quotient therefore has orientation
\(S^3_y\times S^2_z\), equivalently \(S^2_z\times S^3_y\).

\item Let \(t\) be the relative variation of the actual exterior return in
this product. The integral circle Gysin argument can be made
explicit. Put \(P=Y\setminus\operatorname{int}\nu F\), the orbit
complement before its conifold cap. On its physical mapping torus,
whose fourth cohomology vanishes, let \(\eta\)
be a circle connection with \(d\eta=x\), and take
\(d\beta=x^2\). The integral degree-three generator of the circle
total space is represented by
\(-\eta x+\pi^*\beta\); fibre integration and the meridian-one
condition fix its sign. Its boundary restriction in the complex
curve orientation is \(s+k_{\rm Jin}h_F\theta\):
Stokes gives \(k_{\rm phys}=-\int_{F\times S^1}\beta\), and the
physical mapping torus is the inverse of Jin's convention.
Since \(h=-h_F\), its restriction in (K23.2) is
\[
                         t=-k_{\rm Jin}.             \tag{K23.8}
\]
Equivalently, the boundary of a relative \(D^3\times[0,1]\) is
\(S^2\times[0,1]-D^3\times\{1\}+D^3\times\{0\}\), giving precisely
this variation coefficient. The integral Gysin sequence supplies
the integral class; choosing differential representatives does not
replace the integral assertion.

\item The generator \(L\) has variation \(+1\) and quotient
\((y,z)\mapsto(y,\operatorname{Ad}_{\zeta(y)}z)\). In Jin's
identification \((v,0)\sim(g(y)v,1)\), a connection is
\(A_\theta=-\theta\,dg\,g^{-1}\). With \(\omega=dg\,g^{-1}\),
\[
 F=-d\theta\,\omega+(\theta^2-\theta)\omega^2,\qquad
 \int_{S^3}\operatorname{tr}(\omega^3)
                      =-24\pi^2\deg g.
\]
The orientation is \(S^3\times S^1\), so direct integration gives
\(c_2=\deg g\) and \(p_1(\operatorname{ad}g)=-4\deg g\).
The vertical Euler class of this sphere bundle is \(2x\), and
the base \(S^3\times S^1\) has \(H^2=0\). Therefore
\[
                       (k_L,\ell_L)=(-1,-4).
\]
Composing the actual equivariant return with \(L^{-t}\) kills its
variation and gives
\[
                      k_0=0,\qquad \ell_0=\ell-4k.    \tag{K23.9}
\]
Additivity follows from the mapping-torus pair-of-pants bordism
carrying the actual Euler line.

\item Put \({\cal B}=LR^{-1}\). The relative class after (K23.9) is
\({\cal B}^b{\cal C}^s\). Its double is the actual principal circle
bundle over the closed conifold mapping torus. Its tangent bundle
is the pullback of the quotient tangent bundle plus a trivial
line. Since \(H^2=\mathbb Zx\), \(H^4=\mathbb Z[y\theta]\), and
\(x^2=0\), the integral Gysin sequence makes the pullback on \(H^4\)
injective.

\item Represent \({\cal C}^s\) on a disc disjoint from
\(S^3_y\times\{q=1\}\). The representative
\({\cal B}^b{\cal C}^s\) fixes that sphere pointwise and has
normal clutch \(\operatorname{Ad}_{\zeta(y)^b}\) on
\(S^3_y\times S^1\). Its Pontryagin number is \(-4b\), by the
preceding connection calculation. The pullback of \([y\theta]\)
evaluates \(1\) on this cycle. This retains the actual marking:
an isotopy of representatives is the identity on an initial fibre,
and the possible \(x\theta\) ambiguity in a quotient degree-three
lift pulls back to zero. Thus
\[
 \ell-4k=-4b,\qquad b=6A,\qquad n=-6A.              \tag{K23.10}
\]
Here \(n\) is the \(R\)-exponent after rewriting the original
relative class in \(L,R\). This calculation uses the actual Hopf
quotient before the flop; it makes no corresponding quotient claim
after the fixed curve has become two isolated points.

\par\medskip\noindent\textbf{The sign of the central sphere by an oriented plumbing.}\enspace
Let \(\Sigma_f\) denote \(D^4_s\times S^3_t\) glued to
\(S^3_s\times D^4_t\), right coordinate equal to \(f\) of the left
coordinate, with the orientations from
\(\partial(D^4_s\times D^4_t)\). Take three products
\(D_i=D^4_s\times D^4_t\), with orientations \(+,-,+\).
Glue their faces by
\[
 (s_1,t_1)\longmapsto(s_2,t_2)
       =(s_1,s_1t_1s_1^{-1}),\qquad
 (s_2,t_2)\longmapsto(s_3,t_3)
       =(t_2^{-1}s_2t_2,t_2).
\]
The first is on the \(S^3_s\times D^4_t\) faces, the second on the
\(D^4_s\times S^3_t\) faces. Their coordinate orientations agree,
while the adjacent eight-ball signs differ, as required for oriented
gluing. After rounding, the two remaining faces have boundary word
\({\cal A}^{-1}{\cal B}\). The map \({\cal A}\) extends the left
face and \({\cal B}^{-1}\) the right face, so the oriented boundary is
\[
 \Sigma_{{\cal B}^{-1}{\cal A}^{-1}{\cal B}{\cal A}}
                       =\Sigma_{{\cal C}^{-1}}.
\]

\item Orient the first zero-section sphere by \(ds\) in \(D_1\) and
\(-ds\) in \(D_2\), and the second by \(dt\) in \(D_2\) and
\(-dt\) in \(D_3\). At their unique intersection in \(D_2\) their
concatenated orientation is \(-ds\,dt\), the ambient orientation.
Each normal clutch fixes a real line. Thus
\[
 Q=\begin{pmatrix}0&1\\1&0\end{pmatrix},\qquad \sigma(W)=0.
\]
The normal clutches are \(\operatorname{Ad}_s\) and
\(\operatorname{Ad}_{t^{-1}}\). Their normal orientations on their
two hemispheres are respectively \(dt\) and \(-ds\).
Quaternionic conjugation reverses the four-plane orientation and
commutes with every adjoint action, so this second normal-frame
orientation introduces no extra Pontryagin sign. The two
Pontryagin evaluations are \((p,-p)\), with \(|p|=4\).
The eight-manifold has only two four-handles besides its zero-handle;
in particular it is spin. Its unimodular intersection form gives
\[
 p_1^2=(p,-p)Q^{-1}(p,-p)^{\mathsf T}=-32.
\]
The spin Eells--Kuiper formula consequently gives
\[
 \mu(\Sigma_{{\cal C}^{-1}})=-\frac1{28},\qquad
 \mu(\Sigma_{\cal C})=+\frac1{28}.                   \tag{K23.11}
\]

\item There are two explicit reversals in comparing this plumbing with
P3F.22. Its positive Hopf sphere \(u\) has orientation \(-q\),
by (K23.7), and its ordered factors are \((u,y)\), whereas the
plumbing factors are \((y,q)\). Thus its pieces identify as
\[
 D^4_u\times S^3_y=-(S^3_y\times D^4_q),\qquad
 S^3_u\times D^4_y=-(D^4_y\times S^3_q).
\]
The pieces exchange and the ambient orientation reverses. The
literal gluing map is inverted, giving
\(\lambda(F)=-\Sigma_{F^{-1}}\). For a disc-supported class,
putting its supporting disc in the common boundary and cancelling
the opposite product collars gives
\(\Sigma_{F^{-1}}=-\Sigma_F\). Therefore
\(\lambda({\cal C})=\Sigma_{\cal C}\), with the positive sign
in (K23.11) in the run's convention.

\item Finally \(L\) and \(R\) extend the seven-disc, so the central exponent
of the original relative class is its actual class \(A\bmod28\).
Equation (K23.6) adds \(n\) to that exponent, and (K23.10) gives
\(n=-6A\). This proves (K23.1) and the displayed examples.

\item All local orientation matrices and the Chern--Weil integral are
checked exactly in \path{code/part3_fourfold/fixed_edge_signs.py}. The
central word, hyperbolic pairing and rational Eells--Kuiper
arithmetic are checked in \path{code/part3_fourfold/commutator_orientation.py}.
These programs certify those algebraic steps; the actual family,
circle, collar and oriented gluing maps are the constructions in
this proof.

\end{enumerate}
\end{proof}